\documentclass[10pt,a4paper]{amsart}
\usepackage[margin=19mm]{geometry}
\usepackage{amsmath,amssymb,mathtools}
\usepackage[scr=rsfs,cal=euler]{mathalfa}
\usepackage{xfrac}
\usepackage[unicode,hidelinks,bookmarks=false]{hyperref}
\newcommand{\ie}{i.e.\ }
\newcommand{\cf}{cf.\ }
\newcommand{\resp}{respectively\ }
\newcommand{\Subsection}{\subsection}
\providecommand{\nobreakdash}{\nobreak\hbox{-}}
\setlength{\emergencystretch}{2em}
\multlinegap0pt
\usepackage{bm}
\usepackage{amssymb}
\newtheorem{lemma}{Lemma}
\title{Solution to an open problem on the computational complexity of immanant}
\author{Xiangshuai Dong$^a$, Tingzeng Wu$^{b,c}$, Xing Gao$^{d,e}$}
\date{Source: arXiv:2608.24045v1 (CC BY 4.0)}
\begin{document}
\maketitle

\noindent\emph{Source and licence.} Solution to an open problem on the computational complexity of immanant. Version arXiv:2608.24045v1 (CC BY 4.0), distributed by arXiv under the Creative Commons Attribution 4.0 International licence, http://creativecommons.org/licenses/by/4.0/, as recorded in the arXiv licence metadata for this exact version and shown on the abstract page https://arxiv.org/abs/2608.24045v1. Full licence notice: \texttt{licenses/arxiv-2608.24045-CC-BY-4.0.txt}.\par\smallskip

\noindent\textit{Editorial scope.} Counted unit: the article's lemma lem2.6 (source0.tex lines 558--565). The article's complete original proof of it is delivered with it (source0.tex lines 567--719, one \texttt{proof} environment of 153 lines). No mathematics is written, repaired or re-derived: the statement and the proof are the article's own lines, in the article's own notation, and no retained line is substituted or re-flowed.  No further source-local context is needed: every cross-reference of every retained line resolves inside the two delivered spans.  Nothing was omitted from any retained span, so the package carries no omission record.  No retained line contains a citation, so the wrapper carries no bibliography and no bibliography entry.  No retained line contains a figure environment, an image-inclusion command or drawing code, so no image is delivered and the package declares no asset dependency.  Editorial formatting only, in addition: an AMS \texttt{amsart} preamble that restates the article's own declarations and macros used by the retained lines, namely the environments \texttt{lemma} on separate counters, together with the standard packages those lines need.  The retained environments print in this document as Lemma 1 in order of appearance, whereas the article prints the same items with the numbering of its own full document.  The complete native source of the article as distributed in the arXiv e-print for this exact version is bundled unchanged as \texttt{sources/208444/source0.tex}.
\begin{lemma}\label{lem2.6}
Let $n$ be a positive integer, $A=[a_{ij}]$ an $n\times n$ generic variable matrix (i.e., the $a_{ij}$ are independent indeterminates), and set $m(n)=6n^2+2n$. Then there exists an $m(n)\times m(n)$ matrix $\widehat M_n(A)$, each of whose entries is either one of the variables $a_{ij}$ or a constant from $\mathbb{F}$, such that
\begin{eqnarray*}
  \operatorname{Per}_{-2,m(n)}\bigl(\widehat M_n(A)\bigr)=\operatorname{per}(A).
\end{eqnarray*}
Consequently,
$(\operatorname{per}_n)_{n\geq 1}\leq_p(\operatorname{Per}_{-2,m})_{m\geq 1}$.
\end{lemma}

\begin{proof}
For an adjacency matrix $M=[m_{uv}]$ and a permutation $\pi$, the edges corresponding to $\pi$ are $u\to\pi(u)$. Since every vertex chooses exactly one outgoing edge and also exactly one incoming edge, these edges form several vertex–disjoint directed cycles, i.e., a directed cycle cover. The number of cycles is exactly $c(\pi)$, and the weight of the corresponding term is
\[
  \alpha^{c(\pi)}\prod_u m_{u\pi(u)}.
\]
In particular, when $\alpha=-2$, each directed cycle in the cover contributes a factor $-2$. Consider the original matrix $M_n(A)$, with edge weights as defined above. For every $e=(i,j)\in[n]^2$ and $r\in\{1,2,3\}$, the vertex $u_{e,r}$ has only one nonzero–weight outgoing edge, namely $u_{e,r}\longrightarrow v_{e,r}$.
In any nonzero cycle cover, every vertex must choose a nonzero–weight outgoing edge; hence these $3n^2$ edges must appear in every nonzero cycle cover; we shall call them mandatory edges. After choosing these mandatory edges, for each six–vertex subgraph we only need to determine the outgoing edges of the vertices $v_{e,r}$: each may point to some $u_{e,s}$, or, when permitted, to the outside of the six–vertex subgraph. Fix the subgraph corresponding to $e=(i,j)$ and adopt the following convention:
\begin{itemize}
\item $1\in I$ means the external incoming edge $L_i\to u_{e,1}$ is chosen,
      $2\in I$ means the external incoming edge $R_j\to u_{e,2}$ is chosen.
\item $1\in O$ means the external outgoing edge $v_{e,1}\to R_j$ is chosen,
      $2\in O$ means the external outgoing edge $v_{e,2}\to L_i$ is chosen.
\end{itemize}
First fix the choice of all external edges in the whole graph. For each six–vertex subgraph $e$, denote the corresponding label sets by $I_e$ and $O_e$. If $s\in I_e$, then $u_{e,s}$ already has an external incoming edge, so it cannot be pointed to by any $v_{e,r}$. If $r\in O_e$, then the outgoing edge of $v_{e,r}$ already leaves the subgraph, so it cannot connect to any $u_{e,s}$. Thus, the number of $v$ vertices that still need an outgoing edge inside the subgraph is $3-|O_e|$, and the number of $u$ vertices that still need an incoming edge is $3-|I_e|$. If $|I_e|\ne|O_e|$, the numbers differ and a cycle cover cannot be completed, so the contribution of this case is $0$.
When $|I_e|=|O_e|$, an internal completion of this six–vertex subgraph is uniquely represented by a bijection
\[
  \eta_e:\{1,2,3\}\setminus O_e
  \longrightarrow\{1,2,3\}\setminus I_e,
\]
where $\eta_e(r)=s$ means we choose the internal edge $v_{e,r}\to u_{e,s}$. Write
\[
  w_e(\eta_e):=\prod_{r\notin O_e}W_{r,\eta_e(r)}.
\]
The mandatory edges together with the internal edges chosen by $\eta_e$ form in this six–vertex subgraph two kinds of vertex–disjoint directed components: those entirely contained in the six–vertex subgraph (directed cycles), and those that start with an external incoming edge and end with an external outgoing edge (directed paths). Let $\ell_e(\eta_e)$ be the number of directed cycles of the first kind.
The correspondence between external incoming and outgoing edges determined by these directed paths does not depend on $\eta_e$. Indeed, when $|I_e|=|O_e|=0$ there are no such paths; when $|I_e|=|O_e|=1$ there is a unique label for each; when $|I_e|=|O_e|=2$ we have $I_e=O_e=\{1,2\}$, and the two paths connect the mandatory edges of label $1$ and $2$ to the external outgoing edges of the same label. Contract each such directed path together with its two external edges into a single directed edge connecting the external vertices, and temporarily delete the directed cycles completely contained inside the six–vertex subgraph. Doing this for all subgraphs yields a directed graph $\Gamma_0$ that depends only on the already fixed external edges, not on the specific choice of the internal completions $\eta_e$.
If the fixed external edges do not give each $L_i$ and $R_j$ exactly one incoming and one outgoing edge, then no cycle cover exists. Otherwise, let $c_0$ be the number of directed cycles in $\Gamma_0$, and let $w_0$ be the product of the weights of all fixed external edges and mandatory edges. For a choice of internal completions $\boldsymbol\eta=(\eta_e)_{e\in[n]^2}$, denote the corresponding cycle cover by $C(\boldsymbol\eta)$. Contracting directed paths does not change the number of cycles; deleting each internal directed cycle reduces the cycle count by exactly $1$. Hence
\begin{eqnarray*}
  c\bigl(C(\boldsymbol\eta)\bigr)
  &=&c_0+\sum_{e\in[n]^2}\ell_e(\eta_e),\\
  \operatorname{wt}\bigl(C(\boldsymbol\eta)\bigr)
  &=&w_0\prod_{e\in[n]^2}w_e(\eta_e).
\end{eqnarray*}
Thus, after fixing the external edges, summing over all internal completions gives
\begin{eqnarray*}
 &&\sum_{\boldsymbol\eta}
 (-2)^{c(C(\boldsymbol\eta))}\operatorname{wt}(C(\boldsymbol\eta))\\
 &=&(-2)^{c_0}w_0
 \prod_{e\in[n]^2}\left(\sum_{\eta_e}
 (-2)^{\ell_e(\eta_e)}w_e(\eta_e)\right).
\end{eqnarray*}
Therefore the local sum for each six–vertex subgraph is indeed an independent factor. Even if the directed paths belong to a global directed cycle spanning several subgraphs, the factor contributed by the remaining cycles is independent of the internal completion of that subgraph.
Fix a six–vertex subgraph $e=(i,j)$, and abbreviate $I=I_e$, $O=O_e$. Every possible internal completion consists of edges $v_{e,r}\to u_{e,s}$ with weights $W_{rs}$; the $-2$ factor in the local sum is counted only according to the number $\ell_e(\eta_e)$ of directed cycles entirely inside this component.

Suppose that $I=O=\varnothing$. Then each $v_{e,r}$ must point inside the six–vertex subgraph to some $u_{e,s}$, and each of the three $u$ vertices receives exactly one incoming edge. The three mandatory edges $u_{e,r}\to v_{e,r}$ are always present and have weight $1$, so in counting we may regard each pair $(u_{e,r},v_{e,r})$ as a single object labelled $r$. Then choosing edges $v_{e,r}\to u_{e,s}$ corresponds bijectively to permutations in $S_3$, and the number of cycles of the permutation equals the number of directed cycles inside this six–vertex subgraph. Hence the total contribution of all internal edge choices is $\operatorname{Per}_{-2,3}(W)=1$. Suppose $I=O=\{1\}$. Then the external incoming and outgoing edges of label $1$ are already chosen, and we only need to determine the internal edges for labels $2,3$. The sum of all such choices contains the factor
\[
  \operatorname{Per}_{-2}
  \begin{pmatrix}
    W_{22}&W_{23}\\
    W_{32}&W_{33}
  \end{pmatrix}
  =\operatorname{Per}_{-2}
  \begin{pmatrix}
    \frac12&1\\[-1mm]
    -\frac12&-\frac12
  \end{pmatrix}=0.
\]
Thus, regardless of how the external edges outside this subgraph are chosen, the contributions from internal completions with $I=O=\{1\}$ cancel. For $I=O=\{2\}$, the external edges of label $2$ are fixed, and the remaining labels are $1,3$; the relevant submatrix is
\[
  \begin{pmatrix}
    W_{11}&W_{13}\\
    W_{31}&W_{33}
  \end{pmatrix}
  =
  \begin{pmatrix}
    -1&1\\
    1&-\frac12
  \end{pmatrix},
\]
whose $(-2)$-permanent is
\[
  4W_{11}W_{33}-2W_{13}W_{31}
  =4(-1)\left(-\frac12\right)-2(1)(1)=0.
\]
Thus when $I=O$ and both sets are singletons, the local contribution is also $0$.

Suppose that $I=\{1\},O=\{2\}$. Then we must choose edges from $v_{e,1},v_{e,3}$ to $u_{e,2},u_{e,3}$, respectively, and each $u$ vertex is targeted exactly once. Hence there are only two possible completions. The first chooses weights $W_{12}$ and $W_{33}$, forming two directed cycles:
\[
  L_i\to u_{e,1}\to v_{e,1}\to u_{e,2}
  \to v_{e,2}\to L_i,
  \qquad
  u_{e,3}\to v_{e,3}\to u_{e,3}.
\]
The second chooses weights $W_{13}$ and $W_{32}$, forming one directed cycle:
\[
  L_i\to u_{e,1}\to v_{e,1}\to u_{e,3}
  \to v_{e,3}\to u_{e,2}\to v_{e,2}\to L_i.
\]
Both cycle structures use the two external edges $L_i\to u_{e,1}$ and $v_{e,2}\to L_i$, so they are independent of edges outside the subgraph. The two completions differ by one in the number of cycles, so the $(-2)^{c(\pi)}$ factors differ by a factor $-2$. Factoring out the common external edge weight and one $-2$ gives the remaining sum
$(-2)W_{12}W_{33}+W_{13}W_{32}=0$.
Similarly, for $I=\{2\},O=\{1\}$, we choose from $v_{e,2},v_{e,3}$ to $u_{e,1},u_{e,3}$. Choosing $W_{21}$ and $W_{33}$ gives two cycles:
\[
  R_j\to u_{e,2}\to v_{e,2}\to u_{e,1}
  \to v_{e,1}\to R_j,
  \qquad
  u_{e,3}\to v_{e,3}\to u_{e,3}.
\]
Choosing $W_{23}$ and $W_{31}$ gives one cycle:
\[
  R_j\to u_{e,2}\to v_{e,2}\to u_{e,3}
  \to v_{e,3}\to u_{e,1}\to v_{e,1}\to R_j.
\]
Again the number of cycles differs by one, and the remaining sum is
$(-2)W_{21}W_{33}+W_{23}W_{31}=0$.
Hence whenever $I$ and $O$ are both singletons with different labels, the local contribution is also $0$.

Suppose that $I=O=\{1,2\}$. Then all four external edges are chosen. The mandatory edges of labels $1,2$ together with the vertices $L_i,R_j$ form the directed cycle
\[
  L_i\to u_{e,1}\to v_{e,1}\to R_j
  \to u_{e,2}\to v_{e,2}\to L_i .
\]
This directed cycle contributes a factor $-2$, and the edge $L_i\to u_{e,1}$ also contributes the variable factor $a_{ij}$. For labels $3$, the only possible choice is $u_{e,3}\to v_{e,3}\to u_{e,3}$, which forms another directed cycle with contribution $(-2)W_{33}=(-2)\left(-\frac12\right)=1$. Thus the total contribution for $I=O=\{1,2\}$ is $(-2)a_{ij}\cdot 1=-2a_{ij}$.
In summary, denote by
\[
  Z_e(I_e,O_e):=\sum_{\eta_e}
  (-2)^{\ell_e(\eta_e)}w_e(\eta_e)
\]
the local factor corresponding to component $e$ in the factorisation above. The computation shows
\begin{eqnarray*}
  Z_e(I_e,O_e)
  &=&\begin{cases}
    1,&(I_e,O_e)=(\varnothing,\varnothing),\\
    1,&(I_e,O_e)=(\{1,2\},\{1,2\}),\\
    0,&\text{otherwise}.
  \end{cases}
\end{eqnarray*}
In the case $I_e=O_e=\{1,2\}$, the local factor is $(-2)W_{33}=1$; the two contracted boundary paths form a directed cycle through $L_i$ and $R_j$, contributing $-2$, and the external edge $L_i\to u_{e,1}$ contributes $a_{ij}$. Together these give the $-2a_{ij}$ computed above. From the local–global factorisation, if for some six–vertex subgraph the pair $(I_e,O_e)$ is not one of the two nonzero cases, then the total contribution of all cycle covers for that fixed external–edge choice is $0$. Hence we only need to consider subgraphs that either use no external edges or use all four external edges.
For each $i$, among the subgraphs $(i,1),\ldots,(i,n)$ exactly one must use all four external edges. Similarly, for each $j$, among $(1,j),\ldots,(n,j)$ exactly one must use all four. Thus these subgraphs uniquely determine a permutation $\sigma\in S_n$, namely they are exactly $\{(i,\sigma(i)):i\in[n]\}$. Fix such a $\sigma$. After contracting the boundary paths inside each chosen subgraph, for each $i\in[n]$ the component $(i,\sigma(i))$ gives a directed cycle passing through $L_i$ and $R_{\sigma(i)}$. These cycles are vertex–disjoint, and there are no other directed cycles in the contracted graph, so $c_0=n$. The product of the fixed external edge weights is
\[
  w_0=\prod_{i=1}^n a_{i,\sigma(i)},
\]
and all local factors are $1$. Therefore the total contribution of all cycle covers corresponding to $\sigma$ is
\[
  (-2)^n\prod_{i=1}^n a_{i,\sigma(i)}.
\]
Summing over all $\sigma\in S_n$ yields
\begin{eqnarray*}
  \operatorname{Per}_{-2,m(n)}(M_n(A))
  &=&(-2)^n\sum_{\sigma\in S_n}\prod_{i=1}^n a_{i,\sigma(i)}\\
  &=&(-2)^n\operatorname{per}(A).
\end{eqnarray*}
Every nonzero cycle cover contains all mandatory edges $u_{e,r}\to v_{e,r}$. To make the construction completely explicit, for each $i\in[n]$ we fix the mandatory edge $u_{(i,i),1}\longrightarrow v_{(i,i),1}$. These $n$ edges are distinct. Change their weights from $1$ to $-1/2$. This does not change the allowed cycle covers nor the number of cycles in any cover; it merely multiplies the total weight of every nonzero term by $\left(-\frac12\right)^n$.
Denote the modified adjacency matrix by $\widehat M_n(A)$. By linearity of the sum,
\[
  \operatorname{Per}_{-2,m(n)}(\widehat M_n(A))
  =
  (-1/2)^n(-2)^n\operatorname{per}(A)
  =
  \operatorname{per}(A).
\]
Finally, the size of the resulting matrix is $m(n)=6n^2+2n$. Each entry is either one of the original variables $a_{ij}$ or a constant from $\mathbb{F}$. The constants that appear are only
$0,\ 1,\ -1,\ \frac12,\ -\frac12$. Therefore this matrix construction fits the definition of a projection. Since $m(n)$ is quadratic in $n$, these substitutions form a $p$-projection, i.e.,
$(\operatorname{per}_n)_{n\geq1}\leq_p\bigl(\operatorname{Per}_{-2,m}\bigr)_{m\geq1}$.
\end{proof}

\end{document}
